\section{Related Work and Special-Issue Positioning}
\label{sec:related_work}

\subsection{DAG Workflows across Heterogeneous Resources}

Scientific and industrial workflows often use DAGs to express task dependencies, intermediate data flows, and opportunities for parallel execution. In cloud and edge settings, the primary challenge is not only selecting a fast processor for each task, but coordinating computation with the movement of intermediate data. Earlier work on resource allocation and workflow scheduling has shown the importance of accounting for heterogeneous execution environments, workload variation, and latency-sensitive applications \cite{wang2021adaptive,ahmad2014workflow,oluyisola2022dag}. The Compute Continuum expands these concerns because the resource pool spans sensing and edge devices, federated fog systems, cloud services, and high-performance computing resources.

The present work concentrates on the federated fog tier, where tight deadlines and unreliable external connectivity make local coordination essential. This choice complements rather than replaces cloud or HPC execution. A continuum-aware workflow can retain time-critical tasks near the data source while using remote resources for nonurgent, large-scale, or globally coordinated stages. The scheduling formulation therefore exposes the information needed for inter-tier orchestration: task execution requirements, data dependencies, resource eligibility, available bandwidth, and deadlines.

\subsection{Resource-Aware and Data-Aware Scheduling}

Federated fog systems extend conventional fog deployments by interconnecting multiple resource clusters that can exchange computation and storage capacity dynamically \cite{fgcs2024}. This federation can improve elasticity and resilience, but it also makes bandwidth and data-transfer costs central to the placement decision. Recent work has explored adaptive resource provisioning, data-intensive workflow optimisation, and heuristic DAG scheduling in fog and cloud environments \cite{ahmad2014workflow,oluyisola2022dag,ammar2022multicloud}. However, many approaches are computation-centric, assume stable communication conditions, or do not jointly model approximate execution, deadline constraints, and inter-fog data movement.

The proposed formulation directly represents task-level data size and link-level bandwidth. It therefore distinguishes a placement that is computationally attractive from one that is also viable once the required input data have been transferred. The randomized relaxation, greedy scheduler, and genetic scheduler carry this same data- and bandwidth-aware logic into scalable alternatives to the exact MILP. This enables comparisons between global optimisation and practical methods under changing deadline budgets, resource counts, workflow sizes, and failure rates.

\subsection{Resilience and Adaptive Workflow Execution}

Resilience is an important requirement for workflow execution outside a stable data-centre setting. A node failure or loss of a communication path can turn an initially feasible placement into one that violates a deadline. Decentralized communication and fault-tolerant execution methods have been proposed for distributed fog environments, but they may introduce state-replication or synchronization overhead \cite{xie2023decentralized,zhang2022faulttolerant,wang2023latencyaware}. This paper does not claim a new recovery protocol. Instead, it evaluates the scheduling alternatives against increasing workflow-node failure rates and identifies the methods that preserve higher reported average accuracy as failure pressure increases. The result is a scheduling-level resilience view that can be combined with concrete fault-detection and recovery mechanisms in a deployed system.

\subsection{FAIR, Provenance, and Reproducible Workflows}

Compute-Continuum workflows must be interoperable not only at runtime but also as scientific research objects. The FAIR principles call for research objects to be findable, accessible, interoperable, and reusable, and explicitly apply to the algorithms, tools, and workflows that produce scientific data \cite{wilkinson2016fair}. FAIR workflow guidance further emphasises versioned source code, clear metadata, sample inputs and outputs, configuration capture, and persistent identifiers as mechanisms for reproducibility and reuse \cite{devisser2023fairworkflows}.

This study contributes a reproducibility-oriented methodology for the scheduling artifact. The workflow description records the DAG, data dependencies, resource model, configuration, and randomisation settings. A run record can then link these inputs to scheduler version, generated placements, performance metrics, and figures. The methodology does not replace security or access-control policy; instead, it identifies the minimum metadata and provenance objects that a deployment should preserve while respecting the data-governance constraints of the application domain.

\subsection{Compute-Continuum Orchestration Systems}
\label{subsec:rw_orchestration}

Orchestration frameworks for the continuum coordinate placement across tiers that differ in
capability, ownership and reliability. SmartORC organises this as a modular set of services---%
application description, resource discovery and monitoring, a matchmaking solver, and object
storage---and supports both hierarchical and fully distributed orchestration~\cite{carlini2023smartorc}.
Oakestra takes the opposite starting point, arguing that cluster orchestrators designed for the
data centre degrade at the edge precisely because they assume reliable, low-latency,
high-bandwidth links, and replacing them with federated three-tier resource management and
delegated scheduling~\cite{bartolomeo2023oakestra}. Within this journal, Light-HIDRA demonstrates
decentralised resource orchestration for Fog-IoT environments~\cite{nunezgomez2024lighthidra}, and
Sedlak et al. enforce service-level objectives across the continuum by having edge devices build
causal models of their own behaviour~\cite{sedlak2024equilibrium}. These systems answer the
question of \emph{how} a placement decision is propagated and enforced; they treat the decision
itself as a matchmaking or policy step. The present work is complementary: it supplies a formal
statement of that decision in which data-transfer cost is a first-class term, and which an
orchestrator of any of the above designs could invoke at its scheduling point.

\subsection{Scientific and HPC Workflow Engines}
\label{subsec:rw_wms}

Scientific workflow management systems have long automated DAG execution across distributed
resources. Pegasus maps an abstract workflow onto concrete resources and manages the resulting data
movement~\cite{deelman2015pegasus}, and has been extended to span the edge-to-cloud continuum by
forming heterogeneous resource pools with HTCondor~\cite{tanaka2022pegasusedge}.
Snakemake~\cite{koster2012snakemake} and Nextflow~\cite{ditommaso2017nextflow} resolve dependency
graphs and delegate execution to cluster schedulers and container runtimes, prioritising
portability and reproducibility. StreamFlow pairs the workflow graph with a declarative description
of a multi-site execution environment in which sites do not share a common data
space~\cite{colonnelli2021streamflow}, while DagOnStar minimises data movement across local, HPC,
container and cloud targets~\cite{sanchezgallegos2021dagonstar}, and has since been extended with
streaming I/O capabilities that reduce execution time by up to a third on data-intensive
pipelines~\cite{perrotta2025streaming}. These are systems contributions: their placement logic is
matchmaking, rank-based, or delegated to a site scheduler, and inter-site transfer cost enters as
an engineering concern rather than as a decision variable. This paper contributes at the
optimisation layer beneath them, and its descriptor separation
(Section~\ref{sec:reproducibility}) is designed so that the resulting placements can be consumed by
an engine of this class.

\subsection{Bandwidth- and Data-Aware DAG Scheduling}
\label{subsec:rw_dag}

The reference point for DAG scheduling on heterogeneous processors remains HEFT, which ranks tasks
by upward rank over a graph whose edges carry communication cost and assigns each task to the
processor giving the earliest finish time~\cite{topcuoglu2002heft}. HEFT establishes the convention
this paper follows---communication cost lives on the DAG edge---but assumes every task executes to
completion, so it has no mechanism for trading output quality against a deadline. Li et al. address
deadline satisfaction in mobile edge computing by jointly deciding dependent-task placement and
on-demand function configuration, reporting a 29\% reduction in completion time over HEFT with
their GenDoc approximation~\cite{li2023dagmec}. Learned schedulers pursue the same objective
differently: GrapheonRL combines a graph neural network over the dependency structure with a
reinforcement-learning placement policy, reaching ILP-quality solutions substantially faster than
an exact solver~\cite{sharma2025grapheonrl}. The present formulation differs from all three in
admitting partial execution: the fraction $l_{ia}$ makes achieved accuracy a decision variable, so
a workflow that cannot meet its deadline at full fidelity degrades gracefully instead of failing.
It differs from the learned approaches in providing a certificate---the exact MILP and its proofs
bound what any heuristic in the hierarchy can achieve on a given instance.

\subsection{Reproducibility and Provenance Frameworks}
\label{subsec:rw_fair}

The FAIR principles were formulated for data~\cite{wilkinson2016fair} and have since been re-read
specifically for computational workflows, which are simultaneously research objects, software, and
method descriptions~\cite{swilkinson2025fairworkflows}. Concrete standards now exist for each
layer: the Common Workflow Language provides a formal, shared description of a workflow independent
of any one engine~\cite{crusoe2022cwl}, and Workflow Run RO-Crate captures the provenance of an
individual execution at several granularities, bundling inputs, outputs and
code~\cite{leo2024workflowruncrate}. Capturing that provenance is not free: ProvLight shows that
naive capture is prohibitive on constrained edge devices and that a lightweight protocol is
required to make it viable across the edge-to-cloud continuum~\cite{rosendo2023provlight}. The
Workflows Community Summit identifies FAIR computational workflows, time-sensitive workflows and
heterogeneous environments among the field's principal open challenges~\cite{workflowssummit2024}.
This paper does not implement any of these standards. Its contribution is to show that the
scheduling model's inputs and outputs are structurally compatible with them, and to specify the
minimum descriptor and provenance content that a continuum deployment would need to preserve---a
mapping set out in Section~\ref{sec:reproducibility}.


\subsection{Positioning of This Work}

The contribution sits at the intersection of data-aware workflow scheduling, federated fog execution, and reproducible eScience methodology. Its foundational component is the retained correctness and linearization material for the optimisation formulation. Its practical component is the comparison of exact, relaxed, greedy, and genetic scheduling approaches across four dimensions relevant to the Compute Continuum. Its methodological component is a FAIR-oriented experiment-artifact specification that supports future reuse, validation, and interoperability. Together, these elements provide a bridge between optimisation research and deployable Compute-Continuum-aware scientific workflows.
